\begin{theorem}[Critical Goldberg--Seymour theorem]
Every critical multigraph $G$ with
$\chi'(G)\ge\Delta(G)+2$ is elementary.
\end{theorem}
\begin{proof}
Let $H=G^{\circ}$ be the incident core of $G$.  By
\noderef{IncidentCoreReduction}, every vertex of $H$ is incident with an
edge, and
\[
 \operatorname{Critical}(H)\Longleftrightarrow
 \operatorname{Critical}(G),\qquad
 \chi'(H)=\chi'(G),\qquad
 \Delta(H)=\Delta(G),\qquad
 t(H)=t(G).                                             \tag{1}
\]
The same reduction says that elementariness of $H$ implies elementariness
of $G$.  Thus $H$ is critical, and the numerical equalities in (1) transfer
the assumed bound to
\[
                         \Delta(H)+2\leq\chi'(H).        \tag{2}
\]

First suppose that $t(H)<11$.  Criticality of $H$, the no-isolated-vertex
conclusion above, (2), and this order bound are exactly the hypotheses of
\noderef{ThmScheide}.  That theorem makes $H$ elementary, so the final
clause of \noderef{IncidentCoreReduction} makes $G$ elementary.

It remains to consider the complementary case, so $11\leq t(H)$.  Put
$k=\chi'(H)-1$.  Inequality (2) implies $\chi'(H)\geq2$; hence natural-number
subtraction is exact and
\[
                 \chi'(H)=k+1,\qquad \Delta(H)+1\leq k. \tag{3}
\]
Apply \noderef{LargeOrderTerminalWitness} to $H$, using its criticality,
(3), and $11\leq t(H)$.  It supplies an exact edge identity $e_0$, a proper
partial $k$-edge-coloring $\varphi$ leaving exactly $e_0$ uncolored, and a
finite set $X\subseteq V(H)$ which contains both endpoints of $e_0$ and is
both elementary and strongly closed under $\varphi$.  This theorem is the
packaged bounded-attainable-rung endpoint of the Tashkinov-series argument:
its conclusion already supplies the exact terminal witness needed here, so
no additional series hypotheses or correspondence certificate are being
introduced at this step.

Now apply \noderef{egraph2} to $H$ and $k$.  Its hypotheses are precisely
criticality of $H$, the two equalities and inequalities in (3), and the fact
that every vertex of $H$ is incident with an edge.  The data
$(e_0,\varphi,X)$ satisfy its condition~(iv).  The equivalence
(iii)$\leftrightarrow$(iv) first gives condition~(iii), then
(ii)$\leftrightarrow$(iii) gives condition~(ii), and finally
(i)$\leftrightarrow$(ii) gives condition~(i).  Consequently $H$ is
elementary.  The elementary-transfer clause of
\noderef{IncidentCoreReduction} therefore proves that $G$ is elementary.
\end{proof}
